% Source physical 717--719. Native dense two-column Chinese adaptation.
% DM prices are historical, not current prices.
\clearpage
\begingroup
\sffamily\fontsize{8}{10.5}\selectfont
\setlength{\parindent}{0pt}\setlength{\parskip}{3pt}
\setlength{\columnsep}{6mm}
\twocolumn[{\sffamily\Large 数学讲义\par}\vspace{3pt}\hrule\vspace{4pt}
{\sffamily 完整目录单张备索.}\par\vspace{8pt}]
\providecommand{\sgaCatalogEntry}[6]{\noindent\textbf{第 #1 卷:} #2#3. #4 页. #5. DM #6\par}
\sgaCatalogEntry{38}{R. Berger, R. Kiehl, E. Kunz, H.-J. Nastold, }{解析几何中的微分学}{IV, 134}{1967}{12,—}
\sgaCatalogEntry{39}{}{概率论讨论班 I}{II, 189}{1967}{14,—}
\sgaCatalogEntry{40}{J. Tits, }{单李群及其表示的表}{VI, 53}{1967}{6,80}
\sgaCatalogEntry{41}{A. Grothendieck, }{局部上同调}{VI, 106}{1967}{10,—}
\sgaCatalogEntry{42}{J. F. Berglund, K. H. Hofmann, }{紧半拓扑半群与弱殆周期函数}{VI, 160}{1967}{12,—}
\sgaCatalogEntry{43}{D. G. Quillen, }{同伦代数}{VI, 157}{1967}{14,—}
\sgaCatalogEntry{44}{K. Urbanik, }{预测理论讲义}{IV, 50}{1967}{5,80}
\sgaCatalogEntry{45}{A. Wilansky, }{泛函分析专题}{VI, 102}{1967}{9,60}
\sgaCatalogEntry{46}{P. E. Conner, }{周期映射讨论班}{IV, 116}{1967}{10,60}
\sgaCatalogEntry{47}{}{美国中西部范畴讨论班报告 I}{IV, 181}{1967}{14,80}
\sgaCatalogEntry{48}{G. de Rham, S. Maumary, M. A. Kervaire, }{挠率与简单同伦型}{IV, 101}{1967}{9,60}
\sgaCatalogEntry{49}{C. Faith, }{内射模与商环讲义}{XVI, 140}{1967}{12,80}
\sgaCatalogEntry{50}{L. Zalcman, }{解析容量与有理逼近}{VI, 155}{1968}{13,20}
\sgaCatalogEntry{51}{}{概率论讨论班 II}{IV, 199}{1968}{14,—}
\sgaCatalogEntry{52}{D. J. Simms, }{李群与量子力学}{IV, 90}{1968}{8,—}
\sgaCatalogEntry{53}{J. Cerf, }{关于三维球面的微分同胚 ($\Gamma_4=0$)}{XII, 133}{1968}{12,—}
\sgaCatalogEntry{54}{G. Shimura, }{自守函数与数论}{VI, 69}{1968}{8,—}
\sgaCatalogEntry{55}{D. Gromoll, W. Klingenberg, W. Meyer, }{整体 Riemann 几何}{VI, 287}{1968}{20,—}
\sgaCatalogEntry{56}{K. Floret, J. Wloka, }{局部凸空间理论引论}{VIII, 194}{1968}{16,—}
\sgaCatalogEntry{57}{F. Hirzebruch, K. H. Mayer, }{$O(n)$-流形、怪球与奇点}{IV, 132}{1968}{10,80}
\sgaCatalogEntry{58}{}{Riemann 曲面的 Kuramochi 边界}{IV, 102}{1968}{9,60}
\sgaCatalogEntry{59}{K. Jänich, }{可微 $G$-流形}{VI, 89}{1968}{8,—}
\sgaCatalogEntry{60}{G. S. Jones 编, }{微分方程与动力系统讨论班}{VI, 106}{1968}{9,60}
\sgaCatalogEntry{61}{}{美国中西部范畴讨论班报告 II}{IV, 91}{1968}{9,60}
\sgaCatalogEntry{62}{Harish-Chandra, }{半单李群上的自守形式}{X, 138}{1968}{14,—}
\sgaCatalogEntry{63}{F. Albrecht, }{控制理论专题}{IV, 65}{1968}{6,80}
\sgaCatalogEntry{64}{H. Berens, }{处理 Banach 空间上逼近过程的插值方法}{VI, 90}{1968}{8,—}
\sgaCatalogEntry{65}{D. Kölzow, }{测度的微分}{XII, 102}{1968}{8,—}
\sgaCatalogEntry{66}{D. Ferus, }{微分几何与微分拓扑中的全绝对曲率}{VI, 85}{1968}{8,—}
\sgaCatalogEntry{67}{F. Kamber, P. Tondeur, }{平坦流形}{IV, 53}{1968}{5,80}
\sgaCatalogEntry{68}{N. Boboc, P. Mustaţă, }{与二阶椭圆型线性微分算子相联系的调和空间}{VI, 95}{1968}{8,60}
\sgaCatalogEntry{69}{H. Bauer 编, }{位势论讨论班}{VI, 180}{1968}{14,80}
\sgaCatalogEntry{70}{M. H. Löb 编, }{逻辑暑期学校文集}{IV, 331}{1968}{20,—}
\sgaCatalogEntry{71}{}{Pierre Lelong 讨论班 (分析), 1967—1968 年}{VI, 190}{1968}{14,—}
\sgaCatalogEntry{72}{J. Barwise 编, }{无穷语言的句法与语义}{IV, 268}{1968}{18,—}
\sgaCatalogEntry{73}{J. F. Cerf, }{有限群作用讲义}{IV, 123}{1968}{10,—}
\sgaCatalogEntry{74}{A. Fröhlich, }{形式群}{IV, 140}{1968}{12,—}
\sgaCatalogEntry{75}{G. Lumer, }{函数代数与 Hardy 空间}{VI, 80}{1968}{8,—}
\sgaCatalogEntry{76}{R. G. Swan, }{代数 $K$ 理论}{IV, 262}{1968}{18,—}
\sgaCatalogEntry{77}{P.-A. Meyer, }{Markov 过程: Martin 边界}{IV, 123}{1968}{10,—}
\sgaCatalogEntry{78}{H. Herrlich, }{拓扑反射与余反射}{XVI, 166}{1968}{12,—}
\sgaCatalogEntry{79}{A. Grothendieck, }{加性余纤维化范畴与相对余切复形}{IV, 167}{1968}{12,—}
\sgaCatalogEntry{80}{B. Eckmann 编, }{三元组与范畴同调理论讨论班}{IV, 398}{1969}{20,—}
\sgaCatalogEntry{81}{J.-P. Eckmann, M. Guenin, }{统计力学中的代数方法}{VI, 131}{1969}{12,—}
\sgaCatalogEntry{82}{J. Wloka, }{基本空间与广义函数}{VIII, 131}{1969}{12,—}
\sgaCatalogEntry{83}{O. Zariski, }{代数曲面理论引论}{IV, 100}{1969}{8,—}
\sgaCatalogEntry{84}{H. Lüneburg, }{有限置换群的传递扩张}{IV, 119}{1969}{10,—}
\sgaCatalogEntry{85}{P. Cartier, D. Foata, }{交换与重排的组合问题}{IV, 88}{1969}{8,—}
\sgaCatalogEntry{86}{P. Hilton 编, }{范畴论、同调论及其应用 I}{VI, 216}{1969}{16,—}
\sgaCatalogEntry{87}{M. Tierney, }{稳定同伦论中的范畴构造}{IV, 65}{1969}{6,—}
\sgaCatalogEntry{88}{}{概率论讨论班 III}{IV, 229}{1969}{18,—}
\sgaCatalogEntry{89}{M. Behara, K. Krickeberg, J. Wolfowitz 编, }{概率论与信息论}{IV, 256}{1969}{18,—}
\sgaCatalogEntry{90}{N. P. Bhatia, O. Hajek, }{局部半动力系统}{II, 157}{1969}{14,—}
\sgaCatalogEntry{91}{N. N. Janenko, }{求解多维数学物理问题的分步法}{VIII, 194}{1969}{16,80}
\sgaCatalogEntry{92}{P. Hilton 编, }{范畴论、同调论及其应用 II}{V, 308}{1969}{20,—}
\sgaCatalogEntry{93}{K. R. Parthasarathy, }{局部紧群上的乘子}{III, 54}{1969}{5,60}
\sgaCatalogEntry{94}{M. Machover, J. Hirschfeld, }{非标准分析讲义}{VI, 79}{1969}{6,—}
\sgaCatalogEntry{95}{A. S. Troelstra, }{直觉主义原理}{II, 111}{1969}{10,—}
\sgaCatalogEntry{96}{H.-B. Brinkmann, D. Puppe, }{Abel 范畴、正合范畴与对应}{V, 141}{1969}{10,—}
\sgaCatalogEntry{97}{S. O. Chase, M. E. Sweedler, }{Hopf 代数与 Galois 理论}{II, 133}{1969}{10,—}
\sgaCatalogEntry{98}{M. Heins, }{Riemann 曲面上的 Hardy 类}{III, 106}{1969}{10,—}
\sgaCatalogEntry{99}{P. Hilton 编, }{范畴论、同调论及其应用 III}{IV, 489}{1969}{24,—}
\sgaCatalogEntry{100}{M. Artin, B. Mazur, }{平展同伦}{II, 196}{1969}{12,—}
\sgaCatalogEntry{101}{G. P. Szegö, G. Treccani, }{多值变换半群}{VI, 177}{1969}{14,—}
\sgaCatalogEntry{102}{F. Stummel, }{Sobolev 空间中的边值问题与特征值问题}{VIII, 386}{1969}{20,—}
\sgaCatalogEntry{103}{C. T. Taam 编, }{现代分析及其应用讲义 I}{VII, 162}{1969}{12,—}
\sgaCatalogEntry{104}{G. H. Pimbley, Jr., }{非线性算子的特征函数分支及其分岔}{II, 128}{1969}{10,—}
\sgaCatalogEntry{105}{R. Larsen, }{乘子问题}{VII, 284}{1969}{18,—}
\sgaCatalogEntry{106}{S. Mac Lane 编, }{美国中西部范畴讨论班报告 III}{III, 247}{1969}{16,—}
\sgaCatalogEntry{107}{A. Peyerimhoff, }{可和性讲义}{III, 111}{1969}{8,—}
\sgaCatalogEntry{108}{R. M. F. Moss, C. B. Thomas 编, }{代数 $K$ 理论及其几何应用}{IV, 86}{1969}{6,—}
\sgaCatalogEntry{109}{J. Ll. Morris 编, }{微分方程数值解会议}{VI, 275}{1969}{18,—}
\sgaCatalogEntry{110}{G. Chartrand, S. F. Kapoor 编, }{图论的诸多面貌}{VIII, 290}{1969}{18,—}
\sgaCatalogEntry{111}{K. H. Mayer, }{示性数之间的关系}{III, 99}{1969}{8,—}
\sgaCatalogEntry{112}{N. N. Moiseev 编, }{最优化方法讨论会}{IV, 293}{1970}{18,—}
\sgaCatalogEntry{113}{R. Wille, }{同余类几何}{III, 99}{1970}{8,—}
\sgaCatalogEntry{114}{H. Jacquet, R. P. Langlands, }{$\mathrm{GL}(2)$ 上的自守形式}{VII, 548}{1970}{24,—}
\sgaCatalogEntry{115}{K. H. Roggenkamp, V. Huber-Dyson, }{序上的格 I}{XIX, 290}{1970}{18,—}
\sgaCatalogEntry{116}{}{Pierre Lelong 讨论班 (分析), 1969 年}{IV, 195}{1970}{14,—}
\sgaCatalogEntry{117}{Y. Meyer, }{Pisot 数、Salem 数与调和分析}{63}{1970}{6,—}
\sgaCatalogEntry{118}{K. E. Aubert, W. Ljunggren 编, }{第 15 届斯堪的纳维亚数学大会文集, Oslo 1968}{IV, 162}{1970}{12,—}
\sgaCatalogEntry{119}{M. Raynaud, }{群概形与齐性空间上的丰沛层}{III, 219}{1970}{14,—}
\sgaCatalogEntry{120}{D. Siefkes, }{Büchi 的单元二阶后继算术}{XII, 130}{1970}{12,—}
\sgaCatalogEntry{121}{H. S. Bear, }{Gleason 部分讲义}{III, 47}{1970}{6,—}
\sgaCatalogEntry{122}{H. Zieschang, E. Vogt, H.-D. Coldewey, }{曲面与平面不连续群}{VIII, 203}{1970}{16,—}
\sgaCatalogEntry{123}{A. V. Jategaonkar, }{左主理想环}{VI, 145}{1970}{12,—}
\sgaCatalogEntry{124}{P. A. Meyer 编, }{概率论讨论班 IV}{IV, 282}{1970}{20,—}
\sgaCatalogEntry{125}{}{自动证明讨论会}{V, 310}{1970}{20,—}
\sgaCatalogEntry{126}{P. Schapira, }{超函数理论}{XI, 157}{1970}{14,—}
\sgaCatalogEntry{127}{I. Stewart, }{李代数}{IV, 97}{1970}{10,—}
\sgaCatalogEntry{128}{M. Takesaki, }{Tomita 的模 Hilbert 代数理论及其应用}{II, 123}{1970}{10,—}
\sgaCatalogEntry{129}{K. H. Hofmann, }{紧半群与 $C^*$-双代数的对偶性}{XII, 142}{1970}{14,—}
\sgaCatalogEntry{130}{F. Lorenz, }{域上的二次型}{II, 77}{1970}{8,—}
\sgaCatalogEntry{131}{A. Borel 等, }{代数群及相关有限群讨论班}{VII, 321}{1970}{22,—}
\sgaCatalogEntry{132}{}{最优化讨论会}{III, 348}{1970}{22,—}
\sgaCatalogEntry{133}{F. Topsøe, }{拓扑与测度}{XIV, 79}{1970}{8,—}
\sgaCatalogEntry{134}{L. Smith, }{Eilenberg--Moore 谱序列讲义}{VII, 142}{1970}{14,—}
\sgaCatalogEntry{135}{W. Stoll, }{到紧复流形的全纯映射的值分布}{II, 267}{1970}{18,—}
\sgaCatalogEntry{136}{M. Karoubi 等, }{Heidelberg--Saarbrücken--Strasbourg $K$ 理论讨论班}{IV, 264}{1970}{18,—}
\sgaCatalogEntry{137}{S. MacLane 编, }{美国中西部范畴讨论班报告 IV}{III, 139}{1970}{12,—}
\sgaCatalogEntry{138}{D. Foata, M. Schützenberger, }{Euler 多项式的几何理论}{V, 94}{1970}{10,—}
\sgaCatalogEntry{139}{A. Badrikian, }{线性随机函数与柱集测度讨论班}{VII, 221}{1970}{18,—}
\sgaCatalogEntry{140}{C. T. Taam 编, }{现代分析及其应用讲义 II}{VI, 119}{1970}{10,—}
\sgaCatalogEntry{141}{G. Jameson, }{有序线性空间}{XV, 194}{1970}{16,—}
\sgaCatalogEntry{142}{K. W. Roggenkamp, }{序上的格 II}{V, 388}{1970}{22,—}
\sgaCatalogEntry{143}{K. W. Gruenberg, }{群论中的上同调专题}{XIV, 275}{1970}{20,—}
\sgaCatalogEntry{144}{J. A. Yorke 编, }{微分方程与动力系统讨论班 II}{VIII, 268}{1970}{20,—}
\sgaCatalogEntry{145}{E. J. Dubuc, }{丰沛范畴论中的 Kan 扩张}{XVI, 173}{1970}{16,—}
\sgaCatalogEntry{146}{A. B. Altman, S. Kleiman, }{Grothendieck 对偶理论引论}{II, 192}{1970}{18,—}
\sgaCatalogEntry{147}{D. E. Dobbs, }{交换环的 Čech 上同调维数}{VI, 176}{1970}{16,—}
\sgaCatalogEntry{148}{R. Azencott, }{局部紧群的 Poisson 空间}{IX, 141}{1970}{14,—}
\sgaCatalogEntry{149}{R. G. Swan, E. G. Evans, }{有限群与序的 $K$ 理论}{IV, 237}{1970}{20,—}
\sgaCatalogEntry{150}{Heyer, }{局部紧群的对偶性}{XIII, 372}{1970}{20,—}
\sgaCatalogEntry{151}{M. Demazure, A. Grothendieck, }{群概形 I (SGA 3)}{XV, 562}{1970}{24,—}
\sgaCatalogEntry{152}{M. Demazure, A. Grothendieck, }{群概形 II (SGA 3)}{IX, 654}{1970}{24,—}
\sgaCatalogEntry{153}{M. Demazure, A. Grothendieck, }{群概形 III (SGA 3)}{VIII, 529}{1970}{24,—}
\sgaCatalogEntry{154}{A. Lascoux, M. Berger, }{紧 Kähler 流形}{VII, 83}{1970}{8,—}
\sgaCatalogEntry{155}{J. Horváth 编, }{多复变 I, Maryland 1970}{IV, 214}{1970}{18,—}
\sgaCatalogEntry{156}{R. Hartshorne, }{代数簇的丰沛子簇}{XIV, 256}{1970}{20,—}
\sgaCatalogEntry{157}{T. tom Dieck, K. H. Kamps, D. Puppe, }{同伦论}{VI, 265}{1970}{20,—}
\sgaCatalogEntry{158}{T. G. Ostrom, }{有限平移平面}{IV, 112}{1970}{10,—}
\sgaCatalogEntry{159}{R. Ansorge, R. Hass, }{线性与非线性初值问题的差分格式的收敛性}{VIII, 145}{1970}{14,—}
\sgaCatalogEntry{160}{L. Sucheston, }{遍历论与概率论论集}{VII, 277}{1970}{20,—}
\sgaCatalogEntry{161}{J. Stasheff, }{从同伦观点看 $H$-空间}{VI, 95}{1970}{10,—}
\sgaCatalogEntry{162}{Harish-Chandra, van Dijk, }{约化 $p$ 进群上的调和分析}{IV, 125}{1970}{12,—}
\sgaCatalogEntry{163}{P. Deligne, }{具有正则奇点的微分方程}{III, 133}{1970}{12,—}
\sgaCatalogEntry{164}{J. P. Ferrier, }{完备代数讨论班}{II, 69}{1970}{8,—}
\sgaCatalogEntry{165}{J. M. Cohen, }{稳定同伦}{V, 194}{1970}{16,—}
\sgaCatalogEntry{166}{A. J. Silberger, }{$p$ 进数域上的 $\mathrm{PGL}_2$: 表示、球函数与 Fourier 分析}{VII, 202}{1970}{18,—}
\sgaCatalogEntry{167}{Lavrentiev, Romanov, Vasiliev, }{微分方程的多维反问题}{V, 59}{1970}{10,—}
\sgaCatalogEntry{168}{F. P. Peterson, }{Steenrod 代数及其应用: 庆祝 N. E. Steenrod 六十寿辰的会议}{VII, 317}{1970}{22,—}
\sgaCatalogEntry{169}{M. Raynaud, }{Hensel 局部环}{V, 129}{1970}{12,—}
\sgaCatalogEntry{170}{C. T. Taam 编, }{现代分析及其应用讲义 III}{VI, 213}{1970}{18,—}
\sgaCatalogEntry{171}{W. M. Fleischman 编, }{集值映射、选择与 $2^X$ 的拓扑性质}{X, 110}{1970}{12,—}
\sgaCatalogEntry{172}{Y.-T. Siu, G. Trautmann, }{间隙层与凝聚解析子层的延拓}{V, 172}{1971}{16,—}
\sgaCatalogEntry{173}{J. N. Mordeson, B. Vinograde, }{任意纯不可分扩域的结构}{IV, 138}{1970}{14,—}
\sgaCatalogEntry{174}{B. Iversen, }{线性行列式及其在代数曲线族的 Picard 概形中的应用}{VI, 69}{1970}{8,—}
\sgaCatalogEntry{175}{M. Brelot, }{位势论中的拓扑与边界}{VI, 176}{1971}{18,—}
\sgaCatalogEntry{176}{H. Popp, }{代数流形的基本群}{IV, 154}{1970}{16,—}
\sgaCatalogEntry{177}{J. Lambek, }{挠理论、加性语义与商环}{VI, 94}{1971}{12,—}
\sgaCatalogEntry{178}{Th. Bröcker, T. tom Dieck, }{配边理论}{XVI, 191}{1970}{18,—}
\sgaCatalogEntry{179}{}{Bourbaki 讨论班, 1968/69 卷, 第 347—363 讲}{IV, 295}{1971}{22,—}
\sgaCatalogEntry{180}{}{Bourbaki 讨论班, 1969/70 卷, 第 364—381 讲}{IV, 310}{1971}{22,—}
\sgaCatalogEntry{181}{F. DeMeyer, E. Ingraham, }{交换环上的可分代数}{V, 157}{1971}{16,—}
\sgaCatalogEntry{182}{L. D. Baumert, }{循环差集}{VI, 166}{1971}{16,—}
\sgaCatalogEntry{183}{P. F. Hsieh, A. W. J. Stoddart 编, }{微分方程解析理论}{VI, 225}{1971}{20,—}
\sgaCatalogEntry{184}{R. M. Brooks 编, }{多复变讨论会, Park City, Utah, 1970}{V, 234}{1971}{20,—}
\sgaCatalogEntry{185}{J. Horváth 编, }{多复变 II, Maryland 1970}{III, 287}{1971}{24,—}
\sgaCatalogEntry{186}{M. Capobianco, J. B. Frechen, M. Krolik 编, }{图论的新趋势}{VI, 219}{1971}{18,—}
\sgaCatalogEntry{187}{H. S. Shapiro, }{逼近论专题}{VIII, 275}{1971}{22,—}
\clearpage
\onecolumn
\endgroup
